\documentclass[10pt,a4paper]{article}

\usepackage{lmodern}
\usepackage[T1]{fontenc}
\usepackage[utf8]{inputenc}
\usepackage[margin=1in]{geometry}
\usepackage{booktabs}
\usepackage{longtable}
\usepackage{array}
\usepackage{listings}
\usepackage{xcolor}
\usepackage[most]{tcolorbox}
\usepackage{microtype}
\usepackage[hidelinks]{hyperref}
\usepackage{parskip}

\title{STRATOS: Installation Guide}
\author{\textit{Taine J. Rossini}}
\date{Version 1.0.0.}

\begin{document}
\maketitle

\section{}

The following paths are used in the \texttt{makefile}, if yours are different, please edit the \texttt{makefile} accordingly. It is assumed that the user has already appended the \texttt{.bashrc} file according to the instructions in the \texttt{gdtk} installation guide.

\begin{verbatim}
    INSTALL_DIR ?= $(DGD)
    MAIN_DIR := $(DGD_REPO)
\end{verbatim}

\section{}

It is recommended to clean the \texttt{gdtk/src/gas} library before building.

\begin{verbatim}
    cd $DGD_REPO/src/gas
    make clean
\end{verbatim}

\section{}

\texttt{stratos} can then be built and installed by running the following commands,

\begin{verbatim}
    cd $DGD_REPO/examples/stratos
    make install
\end{verbatim}

\section{}

The local \texttt{Python} environment also requires installation of some packages. On \texttt{Linux} this can be performed using the following commands,

\begin{verbatim}
    python3 -m venv .venv
    source .venv/bin/activate
    pip install numpy scipy cffi matplotlib
\end{verbatim}

\end{document}
